\begin{helper}[Original free-color full path]
Fix a defective color $\delta$.  Put
$X=V(T)$, $r=\gammaIndex(\delta)$,
$U_\gamma=\{u^1_{\gamma(r)},u^2_{\gamma(r)}\}$, and
$Z=X\setminus U_\gamma$.  There are a color $\beta$, vertices
$v\in Z$ and $z$, and a simple full $(\delta,\beta)$ incidence path
$P$ from $v$ to $z$ such that $\beta$ is missing at $v$, is free for
$(T,e_0,\varphi)$, differs from $\gamma(r)$, and $P$ has precisely the
support of the two-color component through $v$ and ends outside $X$.
All vertices and edges here are the exact identities of the supplied
reduced configuration and replacement data.
\end{helper}
\begin{proof}
By \noderef{ScheideReducedConfiguration}, the supplied reduced
configuration includes all of the data and laws of
\noderef{ScheideStructuralConfiguration}.  In particular, its ordered
layout has $p$ base vertices followed by $s$ two-vertex blocks, with
\[
             T.length+1=p+2s.                                      \tag{1}
\]
The vertex map of a tree-sequence is injective by
\noderef{TreeSequence}, and \noderef{TreeSequenceVertexSet} defines
$X=V(T)$ to be the image of that map.  Thus (1) gives $|X|=p+2s$.
The two vertices in the $r$-th block occupy the distinct positions
$p+2r$ and $p+2r+1$, so they are distinct members of $X$.  Deleting
them therefore gives
\[
                     |Z|=p+2s-2.                                  \tag{2}
\]

The first edge of \noderef{TreeSequence} is $e_0$, and its attachment
law at the first edge says that the two ends of $e_0$ are the vertices
in positions $0$ and $1$.  The inequality $p\ge2$ and injectivity of
the vertex map separate those positions from the deleted positions
$p+2r,p+2r+1$.  Hence both distinct ends of the loopless edge $e_0$
belong to $Z$.

The structural configuration supplies a globally maximum triple.
Unpacking \noderef{GlobalMaximumTriple}, this says in particular both
that $T$ is a Tashkinov tree and that the hypotheses needed below hold.
By \noderef{ScheideGlobalMaximumOrder}, $X$ is elementary.  Unpacking
the definition \noderef{ColoredVertexSetElementary}, whenever
$x,y\in Z\subseteq X$ and $x\ne y$, the sets of colors missing at $x$
and at $y$ are disjoint.

Now apply \noderef{ScheideGlobalMaximumMissingCounts} to each vertex of
$Z$.  Every vertex contributes at least one missing color, while each of
the two ends of $e_0$ contributes at least two.  Choose one missing
color at every vertex of $Z$ and, at each end of $e_0$, choose a second
one different from the first.  The preceding pairwise-disjointness
shows that all $|Z|+2$ chosen colors are distinct.  Using (2), the set
of colors missing somewhere on $Z$ consequently has cardinality at
least
\[
                  |Z|+2=(p+2s-2)+2=p+2s.                           \tag{3}
\]

We next count the set defined by \noderef{TreeUsedColors}.  Its members
are precisely the colors on displayed edges after the distinguished
first edge.  Among the base edges there are only the positions
$2,\ldots,p-1$, so these contribute at most $p-2$ colors.  Every
remaining displayed edge belongs to one of the $s$ blocks, and the two
edges in block $i$ both have the single color $\gamma(i)$ by
\noderef{ScheideStructuralConfiguration}.  The block edges therefore
contribute at most $s$ further colors.  Hence
\[
             |\varphi\langle T\rangle|\le p-2+s.                   \tag{4}
\]
For completeness, each color counted in (4) is also missing at a base
vertex.  For a nonfirst base edge this follows by applying the
Tashkinov-tree condition \noderef{TashkinovTree} obtained from
\noderef{GlobalMaximumTriple}: its color is missing at an earlier base
vertex.  For a block color, the required base witness is an explicit
field of \noderef{ScheideStructuralConfiguration}.  Every base vertex
lies in $Z$, since its position is below $p$ whereas the two deleted
positions are at least $p$.

Finally, \noderef{ScheideDefectiveSetup} makes
$r=\gammaIndex(\delta)$ an element of the $s$-element index type.
Consequently $s\ge1$, and $p\ge2$ gives
\[
                         p-2+s<p+2s.                               \tag{5}
\]
By (3)--(5), some color $\beta$ is missing at a vertex $v\in Z$ but is
not used on the displayed colored edges.  Since $Z\subseteq X$, this
color is missing on $T$; together with non-use, the definition
\noderef{TreeFreeColors} says exactly that $\beta$ is free.  Moreover
$\gamma(r)$ colors the first displayed edge of the $r$-th block and is
therefore used, so $\beta\ne\gamma(r)$.

Apply \noderef{ScheideDefectiveChainPathExists} with the globally
maximum triple, the supplied membership of $\delta$ in the defective
colors, the free color $\beta$, and the facts $v\in Z\subseteq X$ and
$\beta\in\overline\varphi(v)$.  It supplies a vertex $z$ and a walk
$P$ from $v$ to $z$ satisfying
\noderef{ScheideDefectiveChainPath}.  Unpacking that predicate, $P$ is
simple, $z\notin X$, and the support of $P$ is exactly the full
$(\delta,\beta)$ incidence component through $v$, which is precisely
the claimed conclusion.
\end{proof}
